\section{Conditional arithmetic targets and reassembly ledger}
\label{sec:arithmetic-ledger}

The preceding algebra identifies the first places where unequal-residual
Möbius signs can survive.  This section states what an arithmetic estimate
would have to control and how such an estimate would enter the existing
endpoint ledger.

\subsection{The first rectangle target}

For complete keys \(i\ne j\), define the literal oriented rectangle sum
\begin{equation}
\begin{aligned}
\mathfrak X_{i,j}
:={}&
\sum_{m<n}
\one_{\{s_i(m)\ne s_i(n)\}}
\one_{\{s_j(m)=s_j(n)\}}
\\
&\times
\mu(s_i(m))\mu(s_i(n))
\frac{
\overline{\beta_{i,m}}\beta_{i,n}
\overline{\beta_{j,n}}\beta_{j,m}
}{W_mW_n}.
\end{aligned}
\label{eq:literal-rectangle-target}
\end{equation}
Then
\begin{equation}
 \tr(ER)=2\operatorname{Re}\sum_{i\ne j}\mathfrak X_{i,j}.
 \label{eq:rectangle-target-trace}
\end{equation}
The notation in \eqref{eq:literal-rectangle-target} includes only actual
four-corner rectangles.  A zero coefficient at any corner deletes the
term.  The complete key \(i\) contains its side, external block, opposite
row, \(\gamma\), and time; the same is true independently for \(j\).

There are three inequivalent scopes for a possible estimate:
\begin{enumerate}[label=\textup{(\roman*)},leftmargin=3em]
 \item a bound for one fixed \((i,j)\);
 \item a bound after summing all key pairs incident to one fixed row \(m\);
 \item a bound for the complete trace sum.
\end{enumerate}
The first does not imply the second when the key degree grows.  The third
does not imply uniform local star compensation.  A proof must state which
scope it supplies.

\subsection{Higher one-erasure targets}

For \(q\ge1\), put
\begin{equation}
 \mathfrak L_{2q}(X)
 :=
 S_{2q,1}
 =
 2q\tr(ER^{2q-1}).
 \label{eq:linear-erasure-target}
\end{equation}
By \eqref{eq:literal-residual-boundary-walk}, this is a complete weighted
sum of cross-key closed walks, each carrying one residual pair sign.  A
useful L2 estimate would compare
\[
 |\mathfrak L_{2q}(X)|
\]
with its literal absolute majorant, with constants uniform in all physical
blocks and, if \(q=q_X\) grows, uniform in \(q_X\).

Such an estimate cannot be replaced by:
\begin{itemize}
 \item randomness of independent Möbius signs;
 \item an average over \(h\);
 \item a density-one set of row or key pairs;
 \item a fixed-dimensional computation;
 \item the determinant bijection alone; or
 \item cancellation on the pure \(R\) words, whose target signs are already
       frozen.
\end{itemize}

\subsection{A conditional full-moment certificate}

\begin{proposition}[Conditional moment-to-native certificate]
\label{prop:conditional-moment-certificate}
Let \(q=q_X\ge1\).  Suppose an L2 argument for the complete literal carrier
proves
\begin{equation}
 \tr(R+E)^{2q}
 \le \mathfrak M_{2q}(X).
 \label{eq:L2-full-moment-premise}
\end{equation}
Then
\begin{align}
 \|R+E\|
 &\le \mathfrak M_{2q}(X)^{1/(2q)},
 \label{eq:moment-operator-bound}\\
 \kappa_M
 &\le1+\mathfrak M_{2q}(X)^{1/(2q)},
 \label{eq:moment-native-upper}\\
 \alpha_M
 &\ge
 \left(1-\mathfrak M_{2q}(X)^{1/(2q)}\right)_+.
 \label{eq:moment-native-lower}
\end{align}
If
\[
 \mathfrak M_{2q}(X)
 \le X^{2q\nu_H+o(q)},
\]
then
\[
 \|R+E\|\le X^{\nu_H+o(1)}.
\]
\end{proposition}

\begin{proof}
The defect is Hermitian, so
\[
 \|R+E\|^{2q}\le\tr(R+E)^{2q}.
\]
The native Gram is \(I+R+E\); its extremal eigenvalues give the next two
inequalities.
\end{proof}

Premise \eqref{eq:L2-full-moment-premise} concerns the \emph{complete}
moment.  Bounds for \(\mathfrak L_{2q}\) or \(S_{2q,2}\) do not by
themselves supply it.  In particular, the positivity of \(S_{2q,2}\) is a
structural restriction, not an upper estimate of the correct size.

\subsection{Mass, degree, and represented comparison}

Let \(\mathcal D_{\rm ref}\) be a nonnegative quadratic form.  Suppose, in
one explicitly stated uniform or distinguished scope, that the following
three estimates hold for every coefficient vector \(\zeta\) in that same
scope:
\begin{align}
 \kappa_M&\le X^{\nu_M+o(1)},
 \label{eq:effective-native-exponent}\\
 \langle D_M\zeta,\zeta\rangle
 &\le X^{-\sigma_M+o(1)}\mathcal D_{\rm ref}(\zeta),
 \label{eq:missing-mass-exponent}\\
 \|A_X\zeta\|^2
 &\ge X^{-\lambda_{\rm rep}+o(1)}
       \mathcal D_{\rm ref}(\zeta).
 \label{eq:represented-lower-exponent}
\end{align}
Then the optimal missing-native upper estimate gives the universally
valid comparison
\begin{equation}
 \|M_X\zeta\|^2
 \le
 X^{\nu_M+\lambda_{\rm rep}-\sigma_M+o(1)}
 \|A_X\zeta\|^2.
 \label{eq:mass-degree-ratio}
\end{equation}
Thus the familiar zero-cost condition remains
\begin{equation}
 \boxed{\sigma_M>\nu_M+\lambda_{\rm rep}.}
 \label{eq:zero-cost-mass-degree-gate}
\end{equation}
The present paper changes only how \(\nu_M\) might be bounded.  It does not
prove \eqref{eq:effective-native-exponent},
\eqref{eq:missing-mass-exponent}, or
\eqref{eq:represented-lower-exponent}.

\subsection{The shorted reassembly penalty remains separate}

Let \(A\) be the represented map, \(M\) the missing map, and \(B=A+M\).
Writing \(N_A=\Ker A\), the represented comparison must remove
\[
 W_{\rm nuis}=B(N_A)=M(N_A)
\]
from the full output.  TPC--68 expresses the exact shorted Gram as
\begin{equation}
 G_{\rm sh}=G_{\rm dir}-C_N,
 \qquad C_N\succeq0.
 \label{eq:shorted-penalty}
\end{equation}
The defect \(R+E\) studied here is the atomic-normalized
\emph{missing-native} Gram defect.  It is not automatically the direct
represented-normalized defect in \(G_{\rm dir}\).  A separate physical mass
comparison is required to transfer a bound, and a separate estimate
\[
 \|C_N\|=o(1)
\]
is required to make the nuisance shorting penalty negligible.

Consequently, even the hypothetical estimate
\[
 \|R+E\|=o(1)
\]
would not by itself prove the endpoint reassembly theorem.  It would close
one native gate and leave the missing-mass, represented-frame, and nuisance
gates.

\subsection{Endpoint accounting}

All identities proved here cost zero fixed power of \(X\).  An eventual
arithmetic estimate may introduce:
\[
 \lambda_{\rm native},\qquad
 \lambda_{\rm mass},\qquad
 \lambda_{\rm rep},\qquad
 \lambda_{\rm nuis},
\]
or equivalent losses in the upstream notation.  The complete literal loss,
after signs and quantifier scopes have been matched, must remain strictly
below
\[
 \frac1{400}.
\]
Equality is not enough.  A result outside that budget can still be a valid
operator or arithmetic theorem, but it cannot be promoted to the current
endpoint implication.
